\chapter{Câu hỏi thường gặp}
\label{app:faq}

Phụ lục này trả lời các câu hỏi người đọc thường đặt khi xem báo cáo. Câu trả lời ngắn, kèm
trỏ đến mục liên quan.

\paragraph{Vì sao chọn bàn 11×11 mà không phải 9×9 hay 15×15?}
Bàn 11×11 vừa 121 bit, tiện biểu diễn bằng một số nguyên; bàn 15×15 có 225 bit, vượt 128 bit
và làm bitboard phức tạp hơn. Bàn 9×9 nhỏ hơn, ít không gian cho tám cấp phân hoá. 11×11 cũng
gần với Caro Việt Nam. Chi tiết ở Mục~\ref{sec:gt-tro-choi}.

\paragraph{Vì sao dùng luật tự do, không chặn hai đầu?}
Luật tự do làm điều kiện thắng đơn giản: chỉ cần 5 quân liền, không cần xét hai đầu. Điều này
cho phép engine tính mọi thứ bằng vài phép bit và giữ hàm lượng giá đơn giản. Luật chặn hai
đầu sẽ làm mẫu đe doạ phức tạp hơn, nhưng không cần thiết cho mục tiêu so sánh tám cấp. Chi
tiết ở Mục~\ref{sec:gt-luat} và Mục~\ref{sec:gt-de-doa}.

\paragraph{Vì sao MCTS thua cả tác tử luật?}
Vì MCTS rollout ngẫu nhiên không có tri thức chiến thuật và ở 400 mô phỏng cây quá thưa để thấy
chuỗi ép 4 ply. Tác tử luật có sẵn luật thắng/chặn/đòn kép nên trừng phạt ngay. Đây là kết quả
của cấu hình, không phải kết luận ``MCTS kém'' nói chung. Chi tiết ở Mục~\ref{sec:tn-mcts}.

\paragraph{Vì sao L4 cân bằng 50–50 trong tự đấu nhưng thua L2?}
Vì tự đấu L4 gặp chính nó, cả hai bên đều thiếu tri thức chiến thuật như nhau; kết quả cân
bằng. Khi gặp L2, đối thủ có luật thắng/chặn nên trừng phạt lỗi của L4. Tự đấu cân bằng không
có nghĩa tác tử mạnh. Chi tiết ở Mục~\ref{sec:tn-e3} và Mục~\ref{sec:tn-loi}.

\paragraph{Vì sao L8 mạnh hơn L7 nhưng lợi thế đi trước của L8 nhỏ hơn?}
Vì L8 phòng thủ tốt hơn khi cầm O, biến nhiều ván thua thành hoà. Lợi thế X–O đo mức độ
``quyết định'' của cặp đấu, không đo sức mạnh tuyệt đối. L7 thắng nhanh hơn nhưng cũng thua
nhiều hơn; L8 thắng ít quyết định hơn nhưng ít thua hơn. Chi tiết ở Mục~\ref{sec:tn-e3} và
Mục~\ref{sec:tn-l8-selfplay}.

\paragraph{Vì sao không báo cáo Elo?}
Vì có tác tử thắng/thua gần tuyệt đối (L1 0\%, L8 89,8\%), ước lượng Bradley-Terry không hữu
hạn nếu không có điều chuẩn, và giá trị phụ thuộc điều chuẩn. Báo cáo dùng điểm và khoảng
Wilson, là các đại lượng trực tiếp từ dữ liệu. Chi tiết ở Mục~\ref{sec:lt-thong-ke} và
Phụ lục~\ref{app:thong-ke}.

\paragraph{Vì sao chỉ dùng một seed?}
Vì thời gian. Một seed cho phép chạy nhiều cấu hình hơn; nhiều seed sẽ tốt hơn nhưng vượt ngân
sách. Báo cáo ghi rõ giới hạn này và đề xuất lặp seed ở Phụ lục~\ref{app:cho-bo-sung}.

\paragraph{Vì sao L8 chỉ duyệt 6/50 vòng?}
Vì ngưỡng 55\% trên 100 ván là khắt khe; nhiều ứng viên có điểm 50--55\% vẫn có thể tốt hơn ở
một số khía cạnh. Gating thận trọng giữ bản cũ khi nghi ngờ. Đây là lựa chọn thiết kế, không
phải bằng chứng quá trình thất bại. Chi tiết ở Mục~\ref{sec:tn-l8}.

\paragraph{Vì sao value head quan trọng đến vậy?}
Vì PUCT dùng value để đánh giá lá và REINFORCE dùng value làm baseline. Nếu value sai dấu, tìm
kiếm bị dẫn sai và gradient policy bị nhiễu. Lỗi quy ước dấu là bài học lớn nhất của dự án.
Chi tiết ở Mục~\ref{sec:tn-sl}.

\paragraph{Vì sao không dùng dữ liệu ván người?}
Vì mục tiêu là so sánh các phương pháp tự sinh dữ liệu, không phải đạt trình độ người. Dữ liệu
chuyên gia đến từ Alpha-Beta của nhóm, nên mọi cấp đều tự chứa. Điều này cũng tránh vấn đề bản
quyền và quyền riêng tư.

\paragraph{Làm sao kiểm tra một con số trong báo cáo?}
Mở \texttt{config.yaml} của run, lọc \texttt{games.jsonl} theo tác tử, tính $(W+0.5D)/n$. Hoặc
chạy lại script trong \repoFile{report/scripts/}. Hướng dẫn đầy đủ ở Phụ lục~\ref{app:tai-lap}.

\paragraph{Nếu muốn chơi thử thì làm sao?}
Chạy giao diện web:
\begin{lstlisting}[numbers=none]
uv sync --group dev --extra ui --extra nn
uv run uvicorn caro.ui.server:app
\end{lstlisting}
Mở \url{http://127.0.0.1:8000}, chọn cấp, chơi. Chi tiết ở Phụ lục~\ref{app:ui}.

\paragraph{Kết quả có ngoại suy sang Caro 15×15 không?}
Không trực tiếp. Bàn lớn hơn làm hệ số rẽ nhánh và độ dài ván tăng; lợi thế đi trước có thể
mạnh hơn. Kết quả chỉ áp dụng cho 11×11, luật tự do, bộ khai cuộc và ngân sách đã dùng. Muốn
ngoại suy cần chạy lại.